\section{Results}\label{sec:results}

\subsection{Xi'an: matched-budget comparison}\label{sec:res-xian}
Table~\ref{tab:main_xian_full} compares all methods on the Xi'an instance with ten evaluation seeds. The two primary endpoints of the analysis plan are the hypervolume of the search archive and the expected utility of the deployable front, both for the learned policy against NSGA-II.

% Generated by scripts/make_tables.py. Do not edit by hand.
\begin{table}[t]
\centering
\caption{Main results, Xi'an, full network. Mean $\pm$ standard deviation over the seeds in column $n$. Policy metrics use the deployable greedy front; archive metrics use every evaluated line. $\dagger$: single-line method.}
\label{tab:main_xian_full}
\small
\resizebox{\textwidth}{!}{%
\begin{tabular}{lrrrrrr}
\toprule
Method & $n$ & Policy EUM & Policy HV & Archive HV & IGD$^+$ (policy) & Wall-clock (s) \\
\midrule
RL (ours) & 10 & $-$0.166 $\pm$ 0.018 & 0.579 $\pm$ 0.050 & 0.753 $\pm$ 0.052 & 0.102 $\pm$ 0.021 & 1425.9 $\pm$ 50.8 \\
NSGA-II & 10 & $-$0.106 $\pm$ 0.033 & 0.750 $\pm$ 0.121 & 0.773 $\pm$ 0.123 & 0.093 $\pm$ 0.074 & 230.6 $\pm$ 28.7 \\
Scalarised GA (matched) & 10 & $-$0.135 $\pm$ 0.009 & 0.636 $\pm$ 0.029 & 0.697 $\pm$ 0.027 & 0.129 $\pm$ 0.021 & 221.1 $\pm$ 10.2 \\
Scalarised GA (full budget) & 3 & $-$0.086 $\pm$ 0.004 & 0.814 $\pm$ 0.009 & 0.887 $\pm$ 0.016 & 0.062 $\pm$ 0.006 & 3021.1 $\pm$ 62.0 \\
GA of Michailidis et al. (demand)$^\dagger$ & 10 & $-$0.298 $\pm$ 0.034 & 0.281 $\pm$ 0.063 & 0.485 $\pm$ 0.089 & 0.427 $\pm$ 0.065 & 150.6 $\pm$ 12.4 \\
GA of Michailidis et al. (GGI)$^\dagger$ & 10 & $-$0.272 $\pm$ 0.036 & 0.333 $\pm$ 0.069 & 0.507 $\pm$ 0.084 & 0.358 $\pm$ 0.074 & 150.1 $\pm$ 11.0 \\
Tabular Q-learning (demand)$^\dagger$ & 10 & $-$0.382 $\pm$ 0.121 & 0.198 $\pm$ 0.153 & 0.531 $\pm$ 0.062 & 0.540 $\pm$ 0.231 & 209.6 $\pm$ 12.7 \\
Tabular Q-learning (GGI)$^\dagger$ & 10 & $-$0.432 $\pm$ 0.109 & 0.132 $\pm$ 0.113 & 0.512 $\pm$ 0.043 & 0.642 $\pm$ 0.216 & 208.8 $\pm$ 14.6 \\
Greedy & 1 & $-$0.441 & 0.095 & 0.141 & 0.685 & 302.1 \\
Random search & 10 & $-$0.199 $\pm$ 0.014 & 0.465 $\pm$ 0.032 & 0.471 $\pm$ 0.031 & 0.225 $\pm$ 0.027 & 164.2 $\pm$ 8.2 \\
\bottomrule
\end{tabular}
}
\end{table}


\paragraph{Search quality.} The archive of the RL runs contains every line sampled during training. Its hypervolume is \XianRLArchHv\ ($\pm\,\XianRLArchHvSd$), against \XianNSGAArchHv\ ($\pm\,\XianNSGAArchHvSd$) for NSGA-II. The difference is \XianRLvsNSGAArchHvDiff\ (95\% interval \XianRLvsNSGAArchHvDiffLo\ to \XianRLvsNSGAArchHvDiffHi, $A_{12}=\XianRLvsNSGAArchHvVda$, Holm-adjusted $p=\XianRLvsNSGAArchHvPHolm$), so we cannot distinguish the two. The RL archive is better than that of the matched scalarised GA ($p=\XianRLvsGAMatchedArchHvPHolm$) and far better than those of the two upstream baselines, tabular Q-learning and random lines (all $p\le\XianRLvsTabDemandArchHvPHolm$). By this measure the sampling distribution that the policy learns explores the trade-off surface about as well as NSGA-II does.

\paragraph{Deployable solutions.} The picture changes when we look at what each method hands to a planner for the 15 weight vectors. The expected utility of the RL policy is \XianRLEum\ ($\pm\,\XianRLEumSd$), against \XianNSGAEum\ ($\pm\,\XianNSGAEumSd$) for NSGA-II and \XianGAMatchedEum\ for the matched scalarised GA. NSGA-II is better than the RL policy by \XianRLvsNSGAPolicyEumDiff\ in expected utility (95\% interval \XianRLvsNSGAPolicyEumDiffLo\ to \XianRLvsNSGAPolicyEumDiffHi, $A_{12}=\XianRLvsNSGAPolicyEumVda$, $p=\XianRLvsNSGAPolicyEumPHolm$), and the matched scalarised GA is better too ($p=\XianRLvsGAMatchedPolicyEumPHolm$). The RL policy is better than random lines ($p=\XianRLvsRandomPolicyEumPHolm$) and than the single-line baselines. Under the analysis plan this outcome is a loss for the learned policy on the second primary endpoint, and we state it as such.

\paragraph{Reading the baselines.} The genetic algorithm and the tabular agent of \citet{michailidis2026} return one line optimised for a single reward, so they cannot cover a trade-off set and score below random lines on expected utility. This reflects their design for a single objective, not a defect in their code. The greedy baseline, which picks the best one-step move at every station, is worse still. Greedy lines stall under the straight-line constraint, which shows that the one-step features given to the policy are not enough on their own.

\subsection{Trade-offs and fairness}\label{sec:res-tradeoffs}
% Generated by scripts/make_tables.py. Do not edit by hand.
\begin{table}[t]
\centering
\caption{Raw objective values at the three vertex weights and averaged over the 15 weight vectors, Xi'an, full network. Mean $\pm$ standard deviation over seeds; each value is taken from the deployable front point with the matching weight vector. $\dagger$: single-line method, the same line at every setting.}
\label{tab:objectives_xian_full}
\footnotesize
\resizebox{\textwidth}{!}{%
\begin{tabular}{lrrrr}
\toprule
Method & $D_{\mathrm{tot}}$ & $\mathrm{GGI}_2$ & $s_{\min}$ & $A_{\mathrm{new}}$ \\
\midrule
\multicolumn{5}{l}{\textit{Vertex demand, $w=(1,0,0)$}} \\
\midrule
RL (ours) & 0.0368 $\pm$ 0.0027 & 0.0153 $\pm$ 0.0012 & 0.0118 $\pm$ 0.0020 & 0.4008 $\pm$ 0.0390 \\
RL, demand-only reward & 0.0432 $\pm$ 0.0039 & 0.0152 $\pm$ 0.0017 & 0.0077 $\pm$ 0.0018 & 0.2216 $\pm$ 0.0655 \\
RL, no message passing & 0.0197 $\pm$ 0.0000 & 0.0034 $\pm$ 0.0000 & 0.0006 $\pm$ 0.0000 & 0.3220 $\pm$ 0.0000 \\
RL, no scalarised prior & 0.0356 $\pm$ 0.0027 & 0.0156 $\pm$ 0.0011 & 0.0122 $\pm$ 0.0009 & 0.3729 $\pm$ 0.0490 \\
NSGA-II & 0.0443 $\pm$ 0.0026 & 0.0172 $\pm$ 0.0024 & 0.0120 $\pm$ 0.0033 & 0.2660 $\pm$ 0.0451 \\
Scalarised GA (matched) & 0.0389 $\pm$ 0.0025 & 0.0130 $\pm$ 0.0023 & 0.0057 $\pm$ 0.0034 & 0.2072 $\pm$ 0.0634 \\
Scalarised GA (full budget) & 0.0441 $\pm$ 0.0025 & 0.0142 $\pm$ 0.0027 & 0.0056 $\pm$ 0.0026 & 0.2437 $\pm$ 0.0195 \\
GA of Michailidis et al. (demand)$^\dagger$ & 0.0368 $\pm$ 0.0039 & 0.0154 $\pm$ 0.0013 & 0.0103 $\pm$ 0.0029 & 0.2239 $\pm$ 0.0347 \\
GA of Michailidis et al. (GGI)$^\dagger$ & 0.0371 $\pm$ 0.0045 & 0.0157 $\pm$ 0.0014 & 0.0104 $\pm$ 0.0031 & 0.2611 $\pm$ 0.0402 \\
Tabular Q-learning (demand)$^\dagger$ & 0.0287 $\pm$ 0.0078 & 0.0113 $\pm$ 0.0041 & 0.0076 $\pm$ 0.0042 & 0.2325 $\pm$ 0.0949 \\
Tabular Q-learning (GGI)$^\dagger$ & 0.0264 $\pm$ 0.0078 & 0.0102 $\pm$ 0.0034 & 0.0064 $\pm$ 0.0032 & 0.1826 $\pm$ 0.0749 \\
Greedy & 0.0224 & 0.0083 & 0.0061 & 0.0949 \\
Random & 0.0348 $\pm$ 0.0017 & 0.0129 $\pm$ 0.0015 & 0.0084 $\pm$ 0.0019 & 0.3034 $\pm$ 0.0486 \\
\addlinespace
\multicolumn{5}{l}{\textit{Vertex fairness, $w=(0,1,0)$}} \\
\midrule
RL (ours) & 0.0359 $\pm$ 0.0024 & 0.0160 $\pm$ 0.0013 & 0.0128 $\pm$ 0.0015 & 0.3998 $\pm$ 0.0308 \\
RL, demand-only reward & 0.0384 $\pm$ 0.0148 & 0.0137 $\pm$ 0.0055 & 0.0060 $\pm$ 0.0035 & 0.2232 $\pm$ 0.0667 \\
RL, no message passing & 0.0201 $\pm$ 0.0023 & 0.0095 $\pm$ 0.0007 & 0.0083 $\pm$ 0.0007 & 0.0830 $\pm$ 0.0108 \\
RL, no scalarised prior & 0.0356 $\pm$ 0.0027 & 0.0156 $\pm$ 0.0011 & 0.0122 $\pm$ 0.0009 & 0.3729 $\pm$ 0.0490 \\
NSGA-II & 0.0398 $\pm$ 0.0038 & 0.0187 $\pm$ 0.0010 & 0.0145 $\pm$ 0.0013 & 0.2705 $\pm$ 0.0430 \\
Scalarised GA (matched) & 0.0343 $\pm$ 0.0040 & 0.0168 $\pm$ 0.0010 & 0.0139 $\pm$ 0.0020 & 0.2448 $\pm$ 0.0557 \\
Scalarised GA (full budget) & 0.0398 $\pm$ 0.0017 & 0.0197 $\pm$ 0.0004 & 0.0156 $\pm$ 0.0033 & 0.2249 $\pm$ 0.0445 \\
GA of Michailidis et al. (demand)$^\dagger$ & 0.0368 $\pm$ 0.0039 & 0.0154 $\pm$ 0.0013 & 0.0103 $\pm$ 0.0029 & 0.2239 $\pm$ 0.0347 \\
GA of Michailidis et al. (GGI)$^\dagger$ & 0.0371 $\pm$ 0.0045 & 0.0157 $\pm$ 0.0014 & 0.0104 $\pm$ 0.0031 & 0.2611 $\pm$ 0.0402 \\
Tabular Q-learning (demand)$^\dagger$ & 0.0287 $\pm$ 0.0078 & 0.0113 $\pm$ 0.0041 & 0.0076 $\pm$ 0.0042 & 0.2325 $\pm$ 0.0949 \\
Tabular Q-learning (GGI)$^\dagger$ & 0.0264 $\pm$ 0.0078 & 0.0102 $\pm$ 0.0034 & 0.0064 $\pm$ 0.0032 & 0.1826 $\pm$ 0.0749 \\
Greedy & 0.0213 & 0.0086 & 0.0068 & 0.0834 \\
Random & 0.0304 $\pm$ 0.0028 & 0.0156 $\pm$ 0.0008 & 0.0136 $\pm$ 0.0016 & 0.2606 $\pm$ 0.0423 \\
\addlinespace
\multicolumn{5}{l}{\textit{Vertex access, $w=(0,0,1)$}} \\
\midrule
RL (ours) & 0.0336 $\pm$ 0.0032 & 0.0149 $\pm$ 0.0017 & 0.0124 $\pm$ 0.0020 & 0.4388 $\pm$ 0.0283 \\
RL, demand-only reward & 0.0430 $\pm$ 0.0034 & 0.0159 $\pm$ 0.0008 & 0.0087 $\pm$ 0.0034 & 0.2736 $\pm$ 0.0327 \\
RL, no message passing & 0.0020 $\pm$ 0.0003 & 0.0002 $\pm$ 0.0000 & 0.0000 $\pm$ 0.0000 & 0.3480 $\pm$ 0.0005 \\
RL, no scalarised prior & 0.0356 $\pm$ 0.0027 & 0.0156 $\pm$ 0.0011 & 0.0122 $\pm$ 0.0009 & 0.3729 $\pm$ 0.0490 \\
NSGA-II & 0.0244 $\pm$ 0.0051 & 0.0084 $\pm$ 0.0037 & 0.0047 $\pm$ 0.0031 & 0.4872 $\pm$ 0.0239 \\
Scalarised GA (matched) & 0.0141 $\pm$ 0.0063 & 0.0028 $\pm$ 0.0013 & 0.0006 $\pm$ 0.0004 & 0.4361 $\pm$ 0.0400 \\
Scalarised GA (full budget) & 0.0175 $\pm$ 0.0003 & 0.0034 $\pm$ 0.0002 & 0.0007 $\pm$ 0.0001 & 0.4706 $\pm$ 0.0375 \\
GA of Michailidis et al. (demand)$^\dagger$ & 0.0368 $\pm$ 0.0039 & 0.0154 $\pm$ 0.0013 & 0.0103 $\pm$ 0.0029 & 0.2239 $\pm$ 0.0347 \\
GA of Michailidis et al. (GGI)$^\dagger$ & 0.0371 $\pm$ 0.0045 & 0.0157 $\pm$ 0.0014 & 0.0104 $\pm$ 0.0031 & 0.2611 $\pm$ 0.0402 \\
Tabular Q-learning (demand)$^\dagger$ & 0.0287 $\pm$ 0.0078 & 0.0113 $\pm$ 0.0041 & 0.0076 $\pm$ 0.0042 & 0.2325 $\pm$ 0.0949 \\
Tabular Q-learning (GGI)$^\dagger$ & 0.0264 $\pm$ 0.0078 & 0.0102 $\pm$ 0.0034 & 0.0064 $\pm$ 0.0032 & 0.1826 $\pm$ 0.0749 \\
Greedy & 0.0140 & 0.0052 & 0.0019 & 0.1880 \\
Random & 0.0255 $\pm$ 0.0027 & 0.0103 $\pm$ 0.0015 & 0.0078 $\pm$ 0.0021 & 0.4444 $\pm$ 0.0131 \\
\addlinespace
\multicolumn{5}{l}{\textit{Mean over the 15 weight vectors}} \\
\midrule
RL (ours) & 0.0361 $\pm$ 0.0025 & 0.0157 $\pm$ 0.0011 & 0.0127 $\pm$ 0.0013 & 0.4123 $\pm$ 0.0329 \\
RL, demand-only reward & 0.0420 $\pm$ 0.0065 & 0.0154 $\pm$ 0.0013 & 0.0080 $\pm$ 0.0022 & 0.2481 $\pm$ 0.0157 \\
RL, no message passing & 0.0132 $\pm$ 0.0009 & 0.0041 $\pm$ 0.0007 & 0.0026 $\pm$ 0.0007 & 0.2562 $\pm$ 0.0228 \\
RL, no scalarised prior & 0.0356 $\pm$ 0.0027 & 0.0156 $\pm$ 0.0011 & 0.0122 $\pm$ 0.0009 & 0.3729 $\pm$ 0.0490 \\
NSGA-II & 0.0375 $\pm$ 0.0026 & 0.0155 $\pm$ 0.0015 & 0.0111 $\pm$ 0.0013 & 0.3523 $\pm$ 0.0281 \\
Scalarised GA (matched) & 0.0323 $\pm$ 0.0010 & 0.0129 $\pm$ 0.0006 & 0.0089 $\pm$ 0.0007 & 0.3032 $\pm$ 0.0136 \\
Scalarised GA (full budget) & 0.0362 $\pm$ 0.0006 & 0.0150 $\pm$ 0.0005 & 0.0106 $\pm$ 0.0005 & 0.3406 $\pm$ 0.0034 \\
GA of Michailidis et al. (demand)$^\dagger$ & 0.0368 $\pm$ 0.0039 & 0.0154 $\pm$ 0.0013 & 0.0103 $\pm$ 0.0029 & 0.2239 $\pm$ 0.0347 \\
GA of Michailidis et al. (GGI)$^\dagger$ & 0.0371 $\pm$ 0.0045 & 0.0157 $\pm$ 0.0014 & 0.0104 $\pm$ 0.0031 & 0.2611 $\pm$ 0.0402 \\
Tabular Q-learning (demand)$^\dagger$ & 0.0287 $\pm$ 0.0078 & 0.0113 $\pm$ 0.0041 & 0.0076 $\pm$ 0.0042 & 0.2325 $\pm$ 0.0949 \\
Tabular Q-learning (GGI)$^\dagger$ & 0.0264 $\pm$ 0.0078 & 0.0102 $\pm$ 0.0034 & 0.0064 $\pm$ 0.0032 & 0.1826 $\pm$ 0.0749 \\
Greedy & 0.0214 & 0.0085 & 0.0062 & 0.1352 \\
Random & 0.0309 $\pm$ 0.0013 & 0.0132 $\pm$ 0.0005 & 0.0102 $\pm$ 0.0006 & 0.3347 $\pm$ 0.0180 \\
\addlinespace
\bottomrule
\end{tabular}
}
\end{table}

Table~\ref{tab:objectives_xian_full} gives the raw objectives at the three vertex weights and averaged over the 15 deployable lines. NSGA-II produces a real trade-off. Its demand-vertex lines serve \XianNSGAVertexDemandDtot\ of the total flow and reach \XianNSGAVertexDemandAnew\ of the trip-end mass, while its accessibility-vertex lines serve \XianNSGAVertexAccessDtot\ and reach \XianNSGAVertexAccessAnew. The three vertex solutions of the RL policy are almost the same line: demand \XianRLVertexDemandDtot, \XianRLVertexFairnessDtot\ and \XianRLVertexAccessDtot, and accessibility \XianRLVertexDemandAnew, \XianRLVertexFairnessAnew\ and \XianRLVertexAccessAnew. The policy is therefore not responding to the weight vector in the way a multi-objective policy should, although the scalarised prior makes it respond somewhat: its 15 deployable points contain \XianRLLines\ distinct lines on average, against \XianNSGALines\ for NSGA-II. The worst-served group receives between \XianRLVertexDemandSmin\ and \XianRLVertexFairnessSmin\ of its demand under the RL policy, against up to \XianNSGAVertexFairnessSmin\ for the fairness-vertex lines of NSGA-II. No run produced an undefined Gini index (\XianRLNA\ of \XianRLNpoints\ points for RL, \XianNSGANA\ of \XianNSGANpoints\ for NSGA-II).

\subsection{Leave-one-line-out test on Xi'an}\label{sec:res-holdout}
We remove Line~2 from the instance and ask each method for a 27-station extension. The real Line~2 serves \HoldoutRealDtot\ of the total flow in this instance. Table~\ref{tab:holdout} gives the overlap of each generated line with the real one. A cell counts as matched if it lies within one cell (about 1\,km) of the other line.

\begin{table}[t]
\centering
\caption{Recovery of the real Xi'an Line~2 from the instance without it (27 stations, ten seeds). P@1, R@1, F1@1: precision, recall and F1 within one cell. MSD: mean symmetric distance in cells. Haus.: Hausdorff distance in cells. Orient.: angle between principal axes in degrees. $D$ ratio: demand served by the generated line divided by that of the real line. Rows (a) use the demand-vertex solution, rows (b) the deployable line that minimises $g$ at the balanced weight $(1/3,1/3,1/3)$, chosen among the 15 deployable lines, and rows (c) the best-demand line in the search archive. The last column is the percentile of the seed mean among 10\,000 single random lines.}\label{tab:holdout}
\footnotesize
\setlength{\tabcolsep}{3pt}
\resizebox{\textwidth}{!}{% Generated by scripts/run_holdout.py. Do not edit by hand.
% MSD: mean symmetric distance (cells); Haus.: Hausdorff (cells); Orient.: degrees.
% Pct F1@1: percentile of the mean over seeds within the random null (higher F1 is better;
% for MSD, Haus. and Orient. lower is better, and the null percentile is printed in the JSON).
\begin{tabular}{llrrrrrrrrrrr}
\toprule
Method & Point & Seeds & P@1 & R@1 & F1@1 & F1@0 & F1@2 & MSD & Haus. & Orient. & $D$ ratio & Pct F1@1 \\
\midrule
Real Line 2 & reference & -- & 1.00 & 1.00 & 1.00 & 1.00 & 1.00 & 0.00 & 0.00 & 0.00 & 1.00 & -- \\
\multicolumn{13}{l}{Real Line 2: demand $D_{\mathrm{tot}}=0.0292$, $f_{\mathrm{norm}}=(0.706, 0.573, 0.789)$} \\
\midrule
RL (full) & (a) $w{=}(1,0,0)$ & 10 & 0.24 & 0.25 & 0.24 & 0.16 & 0.33 & 4.05 & 9.5 & 48.7 & 1.11 & 95.9 \\
 & (b) centroid & 10 & 0.23 & 0.24 & 0.23 & 0.13 & 0.32 & 4.21 & 9.9 & 49.5 & 1.08 & 94.2 \\
 & (c) archive best $D$ & 10 & 0.26 & 0.27 & 0.26 & 0.16 & 0.36 & 3.94 & 9.9 & 48.9 & 1.25 & 97.2 \\
\midrule
NSGA-II & (a) $w{=}(1,0,0)$ & 10 & 0.41 & 0.40 & 0.40 & 0.29 & 0.50 & 3.00 & 7.8 & 37.4 & 1.34 & 99.9 \\
 & (b) centroid & 10 & 0.41 & 0.41 & 0.41 & 0.27 & 0.51 & 2.73 & 7.0 & 34.8 & 1.20 & 99.9 \\
 & (c) archive best $D$ & 10 & 0.41 & 0.40 & 0.40 & 0.29 & 0.50 & 3.00 & 7.8 & 37.4 & 1.34 & 99.9 \\
\midrule
Matched scalar GA & (a) $w{=}(1,0,0)$ & 10 & 0.27 & 0.27 & 0.27 & 0.14 & 0.40 & 3.87 & 9.5 & 49.1 & 1.18 & 97.4 \\
 & (b) centroid & 10 & 0.57 & 0.57 & 0.57 & 0.44 & 0.64 & 2.02 & 5.9 & 24.8 & 1.12 & 100.0 \\
 & (c) archive best $D$ & 10 & 0.32 & 0.32 & 0.32 & 0.21 & 0.43 & 3.63 & 9.5 & 46.3 & 1.30 & 99.2 \\
\midrule
GA (demand) & (a) $w{=}(1,0,0)$ & 10 & 0.22 & 0.22 & 0.22 & 0.09 & 0.34 & 4.19 & 10.3 & 49.9 & 1.07 & 93.7 \\
 & (b) centroid & 10 & 0.22 & 0.22 & 0.22 & 0.09 & 0.34 & 4.19 & 10.3 & 49.9 & 1.07 & 93.7 \\
 & (c) archive best $D$ & 10 & 0.22 & 0.22 & 0.22 & 0.08 & 0.34 & 4.37 & 10.9 & 50.8 & 1.10 & 92.8 \\
\midrule
GA (GGI2) & (a) $w{=}(1,0,0)$ & 10 & 0.22 & 0.22 & 0.22 & 0.11 & 0.32 & 4.27 & 11.0 & 49.2 & 1.05 & 92.9 \\
 & (b) centroid & 10 & 0.22 & 0.22 & 0.22 & 0.11 & 0.32 & 4.27 & 11.0 & 49.2 & 1.05 & 92.9 \\
 & (c) archive best $D$ & 10 & 0.22 & 0.22 & 0.22 & 0.07 & 0.36 & 4.05 & 10.6 & 46.4 & 1.09 & 92.8 \\
\midrule
Tabular (eff.) & (a) $w{=}(1,0,0)$ & 10 & 0.19 & 0.20 & 0.20 & 0.09 & 0.30 & 4.94 & 11.9 & 55.0 & 0.91 & 87.6 \\
 & (b) centroid & 10 & 0.19 & 0.20 & 0.20 & 0.09 & 0.30 & 4.94 & 11.9 & 55.0 & 0.91 & 87.6 \\
 & (c) archive best $D$ & 10 & 0.20 & 0.21 & 0.21 & 0.09 & 0.30 & 4.60 & 10.9 & 56.1 & 1.13 & 91.1 \\
\midrule
Tabular (GGI2) & (a) $w{=}(1,0,0)$ & 10 & 0.17 & 0.19 & 0.18 & 0.09 & 0.26 & 5.34 & 12.1 & 53.2 & 0.88 & 82.8 \\
 & (b) centroid & 10 & 0.17 & 0.19 & 0.18 & 0.09 & 0.26 & 5.34 & 12.1 & 53.2 & 0.88 & 82.8 \\
 & (c) archive best $D$ & 10 & 0.27 & 0.27 & 0.27 & 0.16 & 0.37 & 4.06 & 10.8 & 50.2 & 1.19 & 97.3 \\
\midrule
Greedy & (a) $w{=}(1,0,0)$ & 1 & 0.11 & 0.15 & 0.13 & 0.04 & 0.23 & 5.07 & 12.0 & 45.6 & 0.22 & 66.7 \\
 & (b) centroid & 1 & 0.11 & 0.15 & 0.13 & 0.04 & 0.23 & 5.07 & 12.0 & 45.6 & 0.22 & 66.7 \\
 & (c) archive best $D$ & 1 & 0.11 & 0.11 & 0.11 & 0.04 & 0.19 & 6.48 & 11.0 & 82.9 & 0.68 & 61.9 \\
\midrule
Random & (a) $w{=}(1,0,0)$ & 10 & 0.17 & 0.19 & 0.18 & 0.07 & 0.31 & 4.34 & 9.7 & 50.1 & 0.99 & 85.2 \\
 & (b) centroid & 10 & 0.19 & 0.21 & 0.20 & 0.09 & 0.32 & 4.31 & 10.1 & 49.6 & 0.93 & 87.9 \\
 & (c) archive best $D$ & 10 & 0.17 & 0.19 & 0.18 & 0.07 & 0.31 & 4.34 & 9.7 & 50.1 & 0.99 & 85.2 \\
\bottomrule
\end{tabular}
}
\end{table}

\paragraph{Reference.} A method that also picks the best of 15 candidate lines by $g$ gains from the selection alone. A random method with this rule reaches an F1 within one cell of \HoldoutRandomVertexFOne\ ($\pm\,\HoldoutRandomVertexFOneSd$) at the demand vertex and \HoldoutRandomBalancedFOne\ ($\pm\,\HoldoutRandomBalancedFOneSd$) at the balanced weight. We compare the methods with this reference, using Mann--Whitney tests with Holm correction over the three compared methods, and not with single random lines. About \HoldoutNullDeadEnd\,\% of the \HoldoutNullN\ single random lines stop before 27 cells, which lowers their recall.

\paragraph{Result.} NSGA-II reaches an F1 of \HoldoutNSGAVertexFOne\ at the demand vertex ($A_{12}=\HoldoutVsRandomNSGAVertexA$ against the random reference, Holm-adjusted $p=\HoldoutVsRandomNSGAVertexPHolm$) and \HoldoutNSGABalancedFOne\ at the balanced weight ($p=\HoldoutVsRandomNSGABalancedPHolm$). The matched scalarised GA reaches \HoldoutGAMatchedBalancedFOne\ at the balanced weight ($p=\HoldoutVsRandomGAMatchedBalancedPHolm$). The RL policy reaches \HoldoutRLVertexFOne\ and \HoldoutRLBalancedFOne\ at the two settings, which we cannot distinguish from the random reference ($p=\HoldoutVsRandomRLVertexPHolm$ and $\HoldoutVsRandomRLBalancedPHolm$). The standard deviations over seeds are large, for example $\pm\,\HoldoutNSGABalancedFOneSd$ for NSGA-II, so these are statements about averages. The lines of these three methods serve more demand than the real Line~2 (balanced-weight ratios of \HoldoutNSGABalancedRatio\ for NSGA-II, \HoldoutGAMatchedBalancedRatio\ for the matched GA and \HoldoutRLBalancedRatio\ for RL), which is expected because the lines are optimised for the same OD matrix.

\paragraph{Interpretation.} Under the objectives used here, evolutionary multi-objective search recovers part of the real alignment more often than chance, and the learned policy does not. The result is consistent with demand-driven planning of the corridor but is not evidence for it. The OD matrix was recorded while both lines were operating, so demand near Line~2 reflects its presence, and Line~2 opened in 2011, before Line~1 in 2013. A real alignment also reflects cost, geology, heritage and land acquisition, none of which the model sees.
